%=====================================================================
% Chapter: Related Work
% Positions the RH-NBV planner against the receding-horizon exploration,
% volumetric information-gain, and targeted/semantic NBV literature.
% Requires: amsmath, booktabs.
%=====================================================================
\chapter{Related Work}
\label{ch:related-work}

Active vision addresses the question of where a robot should place its
sensor next so that perception improves as quickly as possible. The
strategy proposed in this thesis draws on four largely separate lines of
work: (i)~receding-horizon exploration planning, (ii)~receding-horizon
planning for active object reconstruction, (iii)~volumetric
information-gain formulations, and (iv)~targeted, semantics-aware local
next-best-view planning. This chapter reviews each line in turn, states
the objective functions and strategies used by the original works, and
closes by identifying the gap that the proposed planner fills.

%---------------------------------------------------------------------
\section{Next-Best-View and Frontier-Based Exploration}
\label{sec:rw-nbv}

The next-best-view (NBV) problem---selecting the sensor pose that is
expected to reveal the most new information about a scene---dates back to
\textcite{connolly1985}, who iteratively selected the view exposing the
largest unmapped volume. Frontier-based exploration
\parencite{yamauchi1997} instead steers the sensor towards the boundaries
between known free and unknown space, and remains a strong baseline for
volumetric exploration. Both families are \emph{myopic}: each decision
optimises a single next measurement without considering how it
constrains the measurements that follow. Modern volumetric planners
typically maintain a probabilistic occupancy map such as OctoMap
\parencite{hornung2013}, in which each voxel~$x$ carries an occupancy
probability $P_o(x)$ that is updated with every registered depth
measurement; the planners considered below all operate on such a map or
on a dense-grid variant of it.

%---------------------------------------------------------------------
\section{Receding-Horizon Exploration Planning}
\label{sec:rw-rh-exploration}

\textcite{bircher2016} introduced the receding-horizon
``next-best-view'' planner (NBVP) for the autonomous exploration of
unknown 3D spaces with aerial robots, later extended to unified
exploration and surface inspection in \textcite{bircher2018}. At every
planning iteration, NBVP grows a geometric random tree (RRT) of
collision-free vehicle configurations $\xi$ in the currently known free
space. The quality of a tree node $n_k$ is the cumulative unmapped
volume visible along its branch, discounted exponentially by the path
cost:
%
\begin{equation}
  \mathbf{Gain}(n_k) \;=\; \mathbf{Gain}(n_{k-1})
  \;+\; \mathbf{Visible}\bigl(\mathcal{M}, \xi_k\bigr)\,
        e^{-\lambda c\left(\sigma_{k-1}^{k}\right)},
  \label{eq:rw-bircher-gain}
\end{equation}
%
where $\mathcal{M}$ is the occupancy map,
$\mathbf{Visible}(\mathcal{M},\xi_k)$ is the unmapped volume observable
from configuration $\xi_k$, $c(\sigma_{k-1}^{k})$ is the Euclidean cost
of the path segment connecting $\xi_{k-1}$ to $\xi_k$, and $\lambda$
penalises expensive motions. The defining strategy is the
\emph{receding-horizon} execution: although a whole branch is evaluated,
only its \emph{first edge} is executed, after which the tree is rebuilt
on the updated map and the process repeats. This closes the loop between
planning and perception in the manner of receding-horizon (model
predictive) controllers and makes the planner robust to mapping and
tracking errors. However, NBVP reasons purely about unmapped
\emph{volume}: it carries no notion of a target object, no semantics,
and its gain recursion (Eq.~\ref{eq:rw-bircher-gain}) simply accumulates
the visible volume of successive nodes without simulating how earlier
views along the branch would change the map that later views observe.

%---------------------------------------------------------------------
\section{Receding-Horizon Planning for Active Reconstruction}
\label{sec:rw-rh-reconstruction}

\textcite{zhang2022} transferred the receding-horizon idea from
free-flying exploration to \emph{active 3D reconstruction} with a
high-DoF manipulator carrying an eye-in-hand RGB-D camera. Their view
path planner builds an exploration tree by forward sampling directly in
joint space, which guarantees reachability and bounded pose change by
construction. Each voxel~$x$ contributes its occupancy entropy
%
\begin{equation}
  \mathcal{H}(x) \;=\; -P_o(x)\ln P_o(x) \;-\; \bar{P}_o(x)\ln \bar{P}_o(x),
  \qquad \bar{P}_o(x) = 1 - P_o(x),
  \label{eq:rw-zhang-entropy}
\end{equation}
%
and the information gain of a view $v$ averages this entropy over all
voxels $\mathcal{X}$ intersected by the ray set $\mathcal{R}_v$:
%
\begin{equation}
  \mathcal{IG}(v) \;=\; \frac{1}{n}
  \sum_{\forall r \in \mathcal{R}_v}\;\sum_{\forall x \in \mathcal{X}}
  \mathcal{H}(x).
  \label{eq:rw-zhang-ig}
\end{equation}
%
The utility of a path is again accumulated iteratively with an
exponential path-cost weight, $f(\mathcal{M},\xi_k) =
f(\mathcal{M},\xi_{k-1}) + \mathcal{IG}(v_k)\,
e^{\lambda c\left(p_{k-1}^{k}\right)}$, and, to counter the myopia of
greedy NBV, the robot commits to only the first $l$ of the $m$ nodes of
the best branch before replanning. Their experiments over eight
volumetric information-gain metrics showed the multi-step planner
outperforming single-step NBV with hand-designed view spheres. Still,
the method reconstructs \emph{whole objects}: the gain treats every
uncertain voxel equally, there is no semantic prioritisation of a target
region, and---as the authors note themselves---the planner receives no
feedback from the reconstruction quality of the part of the scene that
actually matters for a downstream task.

%---------------------------------------------------------------------
\section{Volumetric Information-Gain Formulations}
\label{sec:rw-ig-metrics}

Which quantity a planner should maximise is a design decision in its own
right. \textcite{delmerico2018} systematically compared volumetric
information (VI) formulations for active object reconstruction. Their
\emph{occlusion-aware} VI weights the entropy of each voxel by the
likelihood that a ray actually reaches it, i.e.\ by the product of the
emptiness probabilities of the voxels traversed before it:
%
\begin{equation}
  P_v(x_n) \;=\; \prod_{i=0}^{n-1} \bar{P}_o(x_i),
  \qquad
  \mathcal{I}_v(x_n) \;=\; P_v(x_n)\, \mathcal{H}(x_n),
  \label{eq:rw-delmerico-oa}
\end{equation}
%
so that voxels hidden behind probably-occupied space contribute little.
They further proposed VIs that concentrate on task-relevant regions,
most notably the \emph{rear side voxel} and \emph{rear side entropy}
formulations, which restrict the summation to unobserved voxels
$\mathcal{S}_o$ expected to lie on the far side of already-observed
surfaces,
%
\begin{equation}
  \mathcal{I}_b(x) =
  \begin{cases}
    1 & x \in \mathcal{S}_o \\
    0 & x \notin \mathcal{S}_o
  \end{cases},
  \qquad
  \mathcal{I}_n(x) \;=\; \mathcal{I}_b(x)\,\mathcal{I}_v(x),
  \label{eq:rw-delmerico-rse}
\end{equation}
%
since a ray incident on the unknown back side of a surface is very
likely to reveal previously unobserved object structure. The candidate
views themselves are drawn from a static view space (a cylinder or
sphere around the object) and scored greedily with a cost-weighted
utility $\,\mathcal{U}_v = (1-\gamma)\,\mathcal{G}_v / \sum_\mathcal{V}
\mathcal{G} - \gamma\, \mathcal{C}_v / \sum_\mathcal{V} \mathcal{C}$.
These formulations supply the theoretical vocabulary---visibility
weighting, rear-side reasoning, gain--cost trade-off---used throughout
this thesis, but the study concerns one-step view selection from a fixed
view space, not multi-step planning.

%---------------------------------------------------------------------
\section{Targeted, Semantics-Aware Local NBV Planning}
\label{sec:rw-gradient-nbv}

Closest to this work, \textcite{burusa2024} formulated NBV planning as a
\emph{local, targeted} perception problem: rather than covering a whole
scene, the camera must quickly improve the perception of a single target
plant node inside a region of interest (ROI), despite heavy occlusion.
Scene knowledge is maintained in a GPU-resident multi-channel voxel grid
storing, per voxel, the occupancy probability $p_o(x)$, a semantic class
label $c_s(x)$ with probability $p_s(x)$, and an ROI indicator. The
utility of a viewpoint $\xi = \{p^c, p^t\}$ (camera position and target
point) is its expected \emph{semantic} information gain, computed by
differentiable ray sampling: along every ray $r$, $N_r$ points are
sampled and each contributes its semantic Shannon entropy weighted by
the accumulated transmittance,
%
\begin{align}
  I_s(r) &= \sum_{i=1}^{N_r} T(i)\, I_s(i),
  &
  T(i) &= \prod_{j=1}^{i-1}\bigl(1 - p_o(j)\bigr),
  \label{eq:rw-burusa-ray}
  \\
  I_s(i) &= -p_s(i)\log_2 p_s(i) - \bigl(1-p_s(i)\bigr)\log_2\bigl(1-p_s(i)\bigr),
  &
  I_s(\xi) &= \sum_{r \in \mathcal{R}} I_s(r).
  \label{eq:rw-burusa-utility}
\end{align}
%
Because the interpolation of grid values is differentiable with respect
to $\xi$, the planner ascends the utility gradient directly,
$\xi_{k+1} = \xi_k + \alpha\, \partial I_s(\xi_k)/\partial \xi_k$,
avoiding candidate sampling altogether. GradientNBV matched the
reconstruction accuracy of a sampling-based semantic NBV planner while
using ten times fewer ray-tracing calls and producing 28\% shorter
trajectories. Two limitations motivate the present work. First, the
planner is \emph{single-step}: each gradient ascent optimises the very
next view only, so sequences of views are never reasoned about jointly.
Second, gradient ascent is inherently local; in strongly occluded,
non-convex utility landscapes it is prone to premature convergence.
\textcite{verlinden2024pso} demonstrated this empirically by replacing
the gradient step with particle swarm optimisation (PSO), in which a
swarm of candidate viewpoints $x_i$ evolves according to
%
\begin{equation}
  v_i(t+1) = w\,v_i(t) + c_1 r_1 \bigl(p_i - x_i(t)\bigr)
                       + c_2 r_2 \bigl(g - x_i(t)\bigr),
  \qquad
  x_i(t+1) = x_i(t) + v_i(t+1),
  \label{eq:rw-pso}
\end{equation}
%
with $p_i$ and $g$ the personal- and global-best positions: on occluded
scenes the population-based search reached roughly 70\% ROI coverage
where low-learning-rate gradient ascent stagnated near 40\%, at the
price of substantially longer trajectories and higher computation. The
tension is thus between locality (smooth, cheap, but myopic and
trap-prone) and global search (robust but expensive and erratic).

%---------------------------------------------------------------------
\section{Discussion and Research Gap}
\label{sec:rw-gap}

Table~\ref{tab:rw-comparison} summarises the reviewed strategies. The
receding-horizon planners \parencite{bircher2016, bircher2018,
zhang2022} are non-myopic but semantically blind: they maximise unmapped
volume or plain occupancy entropy over whole scenes or objects.
Conversely, the targeted planners \parencite{burusa2024,
verlinden2024pso} are semantics-aware and ROI-focused but plan a single
view at a time. Moreover, none of the reviewed methods simulates the
effect of its own planned measurements on the map \emph{within} the
planning horizon: Bircher's gain recursion
(Eq.~\ref{eq:rw-bircher-gain}) and Zhang's accumulated utility both sum
per-node gains on the \emph{current} map, so two nearby views in one
branch are rewarded twice for the same unknown voxels.

\begin{table}[t]
  \centering
  \caption{Comparison of the reviewed viewpoint-planning strategies with
  the proposed planner.}
  \label{tab:rw-comparison}
  \begin{tabular}{lccccc}
    \toprule
    & \multicolumn{1}{c}{Multi-step}
    & \multicolumn{1}{c}{Semantic,}
    & \multicolumn{1}{c}{Belief forward}
    & \multicolumn{1}{c}{Manipulator,}
    & \multicolumn{1}{c}{Escapes local}\\
    Method
    & (non-myopic)
    & ROI-targeted
    & simulation
    & local perception
    & optima \\
    \midrule
    NBVP \parencite{bircher2016,bircher2018} & \checkmark & -- & -- & -- & --\\
    RH--VPP \parencite{zhang2022}            & \checkmark & -- & -- & \checkmark & --\\
    VI metrics \parencite{delmerico2018}     & -- & -- & -- & -- & --\\
    GradientNBV \parencite{burusa2024}       & -- & \checkmark & -- & \checkmark & --\\
    PSO--NBV \parencite{verlinden2024pso}    & -- & \checkmark & -- & \checkmark & \checkmark\\
    \textbf{RH--NBV (ours)}                  & \checkmark & \checkmark & \checkmark & \checkmark & \checkmark\\
    \bottomrule
  \end{tabular}
\end{table}

The planner proposed in this thesis closes this gap by embedding the
transmittance-weighted \emph{semantic} utility of
Eqs.~\ref{eq:rw-burusa-ray}--\ref{eq:rw-burusa-utility} inside a
receding-horizon objective over $H$-step candidate view sequences
$(\xi_1,\dots,\xi_H)$,
%
\begin{equation}
  J \;=\; \sum_{k=0}^{H-1} \gamma^{k}
  \Bigl[\, f_{\mathrm{inc}}\bigl(\mathcal{M}_{\mathrm{pred}}, \xi_k\bigr)
  \;-\; \lambda\, C\bigl(\xi_{k-1}, \xi_k\bigr) \Bigr],
  \label{eq:rw-ours}
\end{equation}
%
where $\gamma \in (0,1]$ is a discount factor prioritising near-term
gains, $C$ is the Euclidean motion cost weighted by $\lambda$, and
$f_{\mathrm{inc}}$ is an \emph{incremental} information gain evaluated
on a predicted map $\mathcal{M}_{\mathrm{pred}}$ that is rolled forward
after every hypothetical view, so that later steps of a sequence are
only credited for information not already gathered by earlier ones. As
in \textcite{bircher2016}, only the first view of the best sequence is
executed before replanning. Candidate sequences combine tangential
orbit steps around the target with occasional random jumps, retaining
the smooth local trajectories of \textcite{burusa2024} without relying
on gradients, and a stagnation-detection mechanism relocates the camera
when ROI coverage plateaus, addressing the local-optimum failure mode
identified by \textcite{verlinden2024pso}. The rear-side intuition of
\textcite{delmerico2018} is available as an optional occlusion-bonus
term that rewards views uncovering occupied yet semantically uncertain
voxels. The result is, to the best of our knowledge, the first planner
that unites non-myopic receding-horizon planning, semantic
target-focused information gain, and in-horizon belief forward
simulation for local perception with a robotic manipulator.
